% !TeX program = xelatex
% Arabic Dyslexia-Friendly Text Simplification Assistant
% Pipeline architecture, data flow, and feasibility analysis
\documentclass[11pt,a4paper]{article}

% ---------- Layout ----------
\usepackage[margin=2.2cm]{geometry}
\usepackage{fontspec}

% ---------- Graphics / diagrams (must load before polyglossia/bidi) ----------
\usepackage{tikz}
\usetikzlibrary{arrows.meta,shapes.geometric,positioning,fit,backgrounds,calc,decorations.pathreplacing,shadows}
\usepackage{pgfgantt}
\usepackage{graphicx}
\usepackage{caption}
\usepackage{subcaption}

% ---------- Tables / lists / boxes ----------
\usepackage{booktabs}
\usepackage{array}
\usepackage{longtable}
\usepackage{enumitem}
\usepackage{tcolorbox}
\tcbuselibrary{skins,breakable}
\usepackage{xcolor}

% ---------- References / misc (hyperref must load before polyglossia/bidi) ----------
\usepackage{hyperref}
\usepackage{microtype}
\usepackage{parskip}

% ---------- Arabic script support (load last: bidi has strict ordering needs) ----------
\usepackage{polyglossia}
\setmainlanguage{english}
\setotherlanguage{arabic}
\newfontfamily\arabicfont[Script=Arabic]{Amiri}

% ---------- Colors ----------
\definecolor{stagecolor}{RGB}{255,229,204}
\definecolor{iocolor}{RGB}{214,240,214}
\definecolor{blockcolor}{RGB}{214,228,246}
\definecolor{cloudcolor}{RGB}{240,224,247}
\definecolor{devicecolor}{RGB}{214,242,240}
\definecolor{prosgreen}{RGB}{39,116,66}
\definecolor{consred}{RGB}{160,39,39}
\definecolor{prosbg}{RGB}{234,247,238}
\definecolor{consbg}{RGB}{252,235,235}
\definecolor{headnavy}{RGB}{22,45,74}

\hypersetup{
  colorlinks=true,
  linkcolor=headnavy,
  citecolor=headnavy,
  urlcolor=headnavy,
  pdftitle={Arabic Dyslexia-Friendly Text Simplification Assistant},
  pdfauthor={Samsung Innovation Campus Capstone Candidate}
}

% ---------- TikZ styles reused across every diagram ----------
\tikzset{
  block/.style={rectangle, draw=black!70, thick, rounded corners, fill=blockcolor,
    minimum width=2.6cm, minimum height=1.05cm, align=center, font=\small, drop shadow},
  stage/.style={rectangle, draw=black!70, thick, rounded corners, fill=stagecolor,
    minimum width=2.8cm, minimum height=1.2cm, align=center, font=\small\bfseries, drop shadow},
  iodata/.style={rectangle, draw=black!60, thick, rounded corners, fill=iocolor,
    minimum width=2.5cm, minimum height=1cm, align=center, font=\small, drop shadow},
  decision/.style={diamond, draw=black!70, thick, fill=yellow!20, align=center,
    font=\scriptsize, inner sep=1pt, aspect=2.2},
  cloudbox/.style={rectangle, draw=violet!60!black, thick, dashed, rounded corners,
    fill=cloudcolor, inner sep=10pt},
  devicebox/.style={rectangle, draw=teal!60!black, thick, dashed, rounded corners,
    fill=devicecolor, inner sep=10pt},
  lbl/.style={font=\scriptsize\itshape, text=black!60, fill=white, inner sep=1pt},
  myarrow/.style={-{Latex[length=2.8mm,width=2mm]}, thick},
  dimarrow/.style={-{Latex[length=2.4mm,width=1.6mm]}, thick, dashed, black!55},
}

\newcommand{\stagenum}[1]{\textcolor{headnavy}{\textbf{#1}}}

% ---------- Pros/Cons box environments ----------
\newtcolorbox{prosbox}[1][]{
  colback=prosbg, colframe=prosgreen, boxrule=0.6pt, arc=2pt,
  title=\textbf{Pros}, fonttitle=\bfseries, coltitle=white,
  colbacktitle=prosgreen, top=2pt, bottom=2pt, left=4pt, right=4pt, #1
}
\newtcolorbox{consbox}[1][]{
  colback=consbg, colframe=consred, boxrule=0.6pt, arc=2pt,
  title=\textbf{Cons}, fonttitle=\bfseries, coltitle=white,
  colbacktitle=consred, top=2pt, bottom=2pt, left=4pt, right=4pt, #1
}
\newtcolorbox{noteb}[1][]{
  colback=blue!4, colframe=headnavy, boxrule=0.5pt, arc=2pt,
  left=6pt, right=6pt, top=4pt, bottom=4pt, #1
}

\newcommand{\proscons}[2]{%
  \begin{minipage}[t]{0.485\textwidth}
    \begin{prosbox}
      \begin{itemize}[leftmargin=1.1em,itemsep=1pt,topsep=1pt]
        #1
      \end{itemize}
    \end{prosbox}
  \end{minipage}\hfill
  \begin{minipage}[t]{0.485\textwidth}
    \begin{consbox}
      \begin{itemize}[leftmargin=1.1em,itemsep=1pt,topsep=1pt]
        #2
      \end{itemize}
    \end{consbox}
  \end{minipage}
}

\title{\vspace{-1.5cm}\textbf{\Large Arabic Dyslexia-Friendly Text Simplification Assistant}\\[4pt]
\large Pipeline Architecture, Data Flow, and Feasibility Analysis}
\author{Samsung Innovation Campus --- Capstone Candidate Idea Brief}
\date{September 2026}

\begin{document}
\maketitle
\vspace{-0.8cm}

\begin{noteb}
\textbf{Purpose of this document.} This is a working technical brief for internal
team review, not a final specification. It lays out the full three-stage pipeline
(\emph{simplify} $\rightarrow$ \emph{diacritize} $\rightarrow$ \emph{read aloud, synced}),
the training and quality-control architecture behind it, and an honest accounting of
the pros, cons, risks, and open decisions --- so the team can evaluate this idea on the
same footing as the other candidate projects before committing week~0 effort to it.
\end{noteb}

\tableofcontents
\thispagestyle{empty}
\newpage

% =====================================================================
\section{Problem Statement}
% =====================================================================

Dyslexic and cognitive-accessibility readers typically have no difficulty
\emph{seeing} text --- the bottleneck is \emph{processing} it: long compound
sentences, low-frequency vocabulary, and dense syntactic structure all raise the
cognitive load of reading past what working memory and phonological decoding can
comfortably absorb. Arabic script adds three specific, compounding difficulties on
top of this general profile that most existing (Latin-script-first) dyslexia tools
do not address at all:

\begin{itemize}[leftmargin=1.4em]
  \item \textbf{Missing diacritics (tashkeel).} Everyday printed and digital Arabic
  omits the vowel marks (\emph{tashkeel}) that disambiguate pronunciation. Fluent
  readers reconstruct the missing vowels from context and lexical knowledge ---
  precisely the kind of rapid contextual inference that is often impaired in
  dyslexia, so the ambiguity that a fluent reader never notices becomes a hard
  blocker for the target user.
  \item \textbf{Agglutinative morphology.} Prefixes (conjunctions, prepositions,
  the definite article) and suffixes (pronouns, gender/number/case markers) attach
  directly to the stem with no orthographic space. One visual ``word'' can pack in
  what would be three to four separate word tokens in English.
  \item \textbf{Connected (cursive) script.} Every letter changes glyph shape
  depending on its position in the word (isolated / initial / medial / final),
  compounding letter-confusion effects that are already a hallmark difficulty in
  dyslexia.
\end{itemize}

\subsection{A concrete illustration}

\begin{center}
\begin{tabular}{@{}p{3.6cm}p{5.4cm}p{5.4cm}@{}}
\toprule
& \textbf{Undiacritized (as normally written)} & \textbf{Diacritized (tashkeel added)} \\
\midrule
Single word &
{\Large\textarabic{مدرسة}} \newline (ambiguous: could plausibly be read as more than one word/pattern without context) &
{\Large\textarabic{مَدْرَسَةٌ}} \newline \emph{madrasatun} --- ``a school'', vowels now fully specified \\
\addlinespace
Agglutination &
{\Large\textarabic{وسيذهبون}} \newline one unbroken visual token &
\textit{wa + sa + yadh-hab + \=una} $\rightarrow$ ``and'' + ``will'' + ``go'' + ``[they, masc.\ pl.]'' --- four English word-equivalents fused into one Arabic string \\
\bottomrule
\end{tabular}
\end{center}

\begin{noteb}
\textbf{Landscape check.} We confirmed via research that no working
dyslexia-friendly Arabic typeface is currently public, and no deployed Arabic
dyslexia-assistance software was found. Related groundwork does exist ---
research on automatic correction of Arabic dyslexic text, and on dyslexia-friendly
Arabic typeface design --- and should be cited in the pitch as prior art. The
claim to make is ``no one has shipped this,'' not ``no one has thought about
this.''
\end{noteb}

% =====================================================================
\section{Solution Overview: the Three-Stage Pipeline}
% =====================================================================

The core idea is a linear, staged pipeline: raw Arabic text is progressively
transformed into a form that is easier for a dyslexic reader to decode, then
delivered through a synchronized audio-visual reading experience. Each stage is
independently swappable and independently gradeable, which matters a lot for both
engineering risk management and for the SIC pitch (each stage can be demoed and
justified on its own).

\begin{figure}[h!]
\centering
\resizebox{\textwidth}{!}{%
\begin{tikzpicture}[node distance=9mm and 7mm]

\node[iodata] (input) {Raw input\\\footnotesize typed / pasted text};
\node[stage, right=20mm of input] (s1) {\stagenum{1}\ Simplify\\\footnotesize fine-tuned AraT5};
\node[stage, right=16mm of s1] (s2) {\stagenum{2}\ Diacritize\\\footnotesize tashkeel restoration};
\node[stage, dashed, right=16mm of s2] (s3) {\stagenum{3}\ Read aloud\\\footnotesize Arabic TTS + word sync\\\footnotesize(optional)};
\node[iodata, right=16mm of s3] (out) {Output\\\footnotesize audio + karaoke-style\\\footnotesize highlighted text};

\draw[myarrow] (input) -- (s1);
\draw[myarrow] (s1) -- (s2) node[midway, lbl, above] {simple text};
\draw[myarrow] (s2) -- (s3) node[midway, lbl, above] {vocalized text};
\draw[myarrow] (s3) -- (out);
\draw[dimarrow] (s2.south) to[out=-90,in=-90,looseness=1.3]
  node[lbl, below]{if Stage~3 dropped} (out.south);

\end{tikzpicture}}
\caption{End-to-end request-time pipeline. Input is typed or pasted text; OCR
is out of scope for this capstone. Stage~3 (dashed) is optional --- if no
adequately open TTS option is found (Section~\ref{sec:tts-arch}), the
pipeline can ship with Stage~2's vocalized text as the final output, dropping
audio and word-sync highlighting entirely.}
\label{fig:pipeline-overview}
\end{figure}

\subsection{Why this stage order, specifically}

\begin{noteb}
\textbf{Simplify before diacritizing, not after.} Diacritization tools generally
assume grammatically well-formed input and are trained on fully-formed sentences;
running diacritization on the \emph{original} complex sentence and then simplifying
a vocalized sentence risks producing diacritic placements that no longer match once
words are reordered, dropped, or replaced. Simplifying first means the diacritizer
only ever has to vocalize sentences the model actually keeps, and word-level
diacritics can be regenerated fresh on the final, simplified surface form.
\end{noteb}

\newpage
% =====================================================================
\section{Stage-by-Stage Detail}
% =====================================================================

\subsection{Stage 1: Simplification (fine-tuned AraT5)}

This is the intellectual core of the project: a sequence-to-sequence model that
takes a complex Arabic sentence and produces a shorter, lower-vocabulary paraphrase
that preserves the original meaning.

\begin{figure}[h!]
\centering
\begin{tikzpicture}[node distance=10mm and 9mm]
\node[iodata] (in) {Complex sentence\\\footnotesize e.g.\ long compound clause,\\\footnotesize low-frequency vocabulary};
\node[stage, right=18mm of in] (model) {Fine-tuned AraT5\\\footnotesize encoder--decoder\\\footnotesize seq2seq};
\node[iodata, right=18mm of model] (out) {Simplified sentence\\\footnotesize shorter, common words,\\\footnotesize same meaning};

\draw[myarrow] (in) -- (model);
\draw[myarrow] (model) -- (out);

\node[block, below=13mm of model, fill=blockcolor] (gate) {Runtime faithfulness gate\\\footnotesize (optional) --- DSPy judge\\\footnotesize scores output live};
\draw[dimarrow] (out.south) |- (gate.east);
\draw[dimarrow] (gate.west) -| node[lbl, near start, left]{fallback if score low} (in.south);
\end{tikzpicture}
\caption{Stage 1 in isolation, including the optional runtime faithfulness gate
described in Section~\ref{sec:qc}.}
\label{fig:stage1}
\end{figure}

\subsubsection{Architecture decision: encoder--decoder, not decoder-only}
\label{sec:why-arat5-decision}

This decision is already settled for the deployed model, but it is worth spelling
out the reasoning in full, since it is the single most defensible technical claim
in the whole pitch.

\begin{center}
\begin{tabular}{@{}p{4.6cm}p{5.6cm}p{5.6cm}@{}}
\toprule
& \textbf{AraT5 (encoder--decoder, seq2seq)} & \textbf{Decoder-only Arabic SLM}\newline\footnotesize(Jais, ALLaM, SILMA, or multilingual Qwen/Gemma) \\
\midrule
Native task fit &
Text simplification is transduction (input $\to$ transformed output, same
meaning) --- the textbook fit for an encoder that fully conditions on the
whole source before the decoder generates anything. &
Trained for open-ended next-token generation; simplification is coerced via
prompting rather than being the native objective. \\
\addlinespace
Faithfulness / groundedness
\newline\footnotesize(the evidence here is weaker than earlier drafts of this
document claimed --- see the note below the table) &
Anchored to the source sequence by cross-attention, but \emph{not} more
faithful in practice: in error-based human evaluation of sentence
simplification, a fine-tuned Control-T5 produced the most errors of any
system tested (350 vs 211 for GPT-4 and 172 for Qwen2.5-72B) and the most
lexical meaning alterations (176 vs 94 and 59) (Wu \& Arase, 2025,
\texttt{doi:10.1145/3744744}). Insertion, deletion and substitution errors
also occur in the reference corpora themselves, undetected by standard
metrics (Devaraj et al., ACL 2022, \texttt{2022.acl-long.506}). &
Prone to ``assistant drift'' --- added explanation, looser rephrasing,
hallucination --- \emph{at deployable sizes}. In the same study the only
small model, Llama-3.2-3B, produced roughly twice the errors of the 72B
model, with Repetition and Hallucination notably more frequent. ChatGPT
over-corrects and does not follow the minimal-edit principle (Fang et al.,
2023, \texttt{arXiv:2304.01746}; Coyne et al., 2023,
\texttt{arXiv:2303.14342}), a tendency that survives fine-tuning: a 770M
T5-large beat fine-tuned 7B/13B LLaMA-2 on precision, 78.0 vs 72.3/74.6
(Park, Do \& Lee, 2025, \texttt{arXiv:2509.20811}). Silently distorting
meaning for a vulnerable reader is a real harm here, not just a quality bug. \\
\addlinespace
Model size / on-device fit &
Far smaller footprint; substantially easier to quantize and run on-device,
directly matching the program's stated on-device AI interest. &
Heavier to deploy at comparable fluency; harder to hit real-time,
low-memory mobile targets. \\
\addlinespace
Out-of-the-box fluency &
Weaker zero-shot; needs fine-tuning to be good at this specific task. &
Stronger out-of-the-box general fluency and world knowledge. \\
\addlinespace
Direct precedent &
AraT5-MSAizer (ACL Anthology \texttt{2024.osact-1.16}) used AraT5 for an
architecturally similar dialect$\to$MSA transformation task. &
No equally direct architectural precedent found for this exact
faithfulness-critical transduction use case. \\
\bottomrule
\end{tabular}
\end{center}

\begin{noteb}
\textbf{Evidence note (2026-09-20).} Earlier revisions of this document argued
that an encoder--decoder is \emph{inherently} more faithful than a decoder-only
model. That claim was never sourced, and the most directly relevant study ---
error-based human evaluation on sentence simplification itself --- contradicts
it: the fine-tuned encoder--decoder made more meaning errors than GPT-4 and
Qwen2.5-72B (Wu \& Arase, 2025). The claim has been narrowed to one the evidence
supports: \emph{at deployable size}, a fine-tuned encoder--decoder remains our
choice, because the only systems that beat it on faithfulness are one to two
orders of magnitude too large to ship on-device, and the one small decoder-only
model tested was no better (326 errors against Control-T5's 350 --- the two
worst systems in the study were precisely the two deployable-size ones).
Among models we could actually ship, the evidence shows \emph{no} faithfulness
advantage for either architecture; the choice rests on task-nativeness,
controllability and size. Architecture is a deployment argument here, not a
faithfulness guarantee; the faithfulness claim rests on Bayan's own human
meaning check.
Full evidence tables: \texttt{references/02\_text\_simplification/architecture\_decisions.md} §1.1.
\end{noteb}

\textbf{Verdict (already settled):} fine-tune AraT5 as the \emph{deployed}
simplification model. A strong API-based Arabic-capable LLM is used only as a
\emph{data-generation} tool (synthetic training pairs) and, separately, as one
candidate for the evaluation judge (Section~\ref{sec:qc}) --- it is never itself
shipped on-device.

\subsubsection{Training data pipeline}
\label{sec:datapipeline}

\begin{noteb}
\textbf{Revised design (this revision).} The original open-ended ``sample
complex BAREC sentences, simplify, filter'' plan has been replaced with a
concrete, level-aware design, arrived at after directly downloading and
profiling the full BAREC corpus (level distribution, \texttt{Domain}/\texttt{Source}/\texttt{Text\_Class}
breakdowns, and the fragment-rate check below) rather than planning against
the paper alone.
\end{noteb}

\paragraph{Level 5 is dropped.}
\label{sec:level5}
BAREC's numeric level~5 (of the 5-level collapse of its native 19 levels) is
only 1,840 of the corpus's 69,441 sentences (2.6\%) --- too small to sample
evenly against the other levels, and by far the largest single semantic jump
if used as a simplification source. Checking what would actually be lost by
dropping it: level~5 holds only 1,307 of the corpus's 18,964
\texttt{Specialized}-tagged sentences (6.9\%) --- the other 93\% of genuinely
hard, domain-specific content is already well represented at levels~3--4
(level~4 alone carries 32.3\% of all \texttt{Specialized} sentences). Dropping
level~5 therefore costs little real coverage of hard content, removes the
tier that would need awkward padding to reach usable volume, and removes the
single riskiest transformation (a direct level-5-to-easy rewrite) from the
generation step entirely.

\paragraph{Two bands, not five levels.} The remaining levels 1--4 (67,601
sentences) collapse into two working bands: \textbf{easy} (levels 1--2,
35,857 sentences) and \textbf{hard} (levels 3--4, 31,744 sentences). Levels~1
and~2, and separately~3 and~4, are hard to reliably tell apart even in BAREC's
own annotation methodology (the 7/5/3-level collapses exist precisely because
the native 19-level boundaries are fuzzy); asking a generator to hit one
specific sub-level rather than a band would bake that fuzziness into the
training signal as noise.

\paragraph{Easy anchors need no generation at all, and are used unfiltered.}
Every easy-band sentence is used directly as an \emph{identity pair}
(input~=~output) --- teaching the model the equally important skill of
leaving already-simple text alone rather than editing it needlessly. This is
free (no LLM call) and, because it uses the whole easy-band pool rather than
a small sample, gives broad coverage of what real simple Arabic actually
looks like.\footnote{An earlier revision of this pipeline filtered the
identity pool for fragments/headers/citations first (a direct check of
BAREC's own level-1 sentences found roughly 59\% look like fragments,
headers, or citations rather than full prose by a simple heuristic; level~2
was cleaner at $\sim$32\%).
That filtering step has been removed on team review: a reader genuinely
encounters headers, dates, and citations while reading, so the model should
see them too --- and because identity pairs are safe by construction
(output is always exactly the input, whatever it is), there is no failure
mode where an unfiltered fragment teaches the model to \emph{generate}
malformed text. The equivalent concern for the hard band doesn't apply
either: the equivalence validator and level classifier
(Section~\ref{sec:qc}) gate the
\emph{generated output}'s quality regardless of whether the anchor was a
full sentence or a fragment, so a separate anchor-quality filter was doing
no protective work that the existing design doesn't already do more
directly. See \texttt{scripts/barec\_simplification\_pipeline.py}.}

\paragraph{Hard anchors are reranked, not single-shot.} Each hard-band
sentence goes through a DSPy generator that proposes 5 candidate
simplifications in one batched call (targeting the easy band, not a
specific sub-level, for the same fuzziness reason as above), a DSPy
validator that scores all 5 candidates' semantic equivalence (with a
reasoning string per candidate) in a second batched call, and a locally
hosted fine-tuned classifier (quantized MARBERT, no API call) that scores
each candidate's readability level; the best-scoring candidate that passes
both checks is kept. Full signatures and the fine-tuned classifier in
Section~\ref{sec:qc}.

\begin{figure}[h!]
\centering
\resizebox{\textwidth}{!}{%
\begin{tikzpicture}[node distance=8mm and 8mm]
\node[iodata] (barec) {BAREC (train+dev+test)\\\footnotesize 69,441 sentences\\\footnotesize levels 1--5};
\node[block, right=14mm of barec] (drop5) {Drop level~5\\\footnotesize 1,840 sentences (2.6\%)};
\node[block, right=14mm of drop5] (sample) {Stratified sample\\\footnotesize by level};

\node[iodata, above right=8mm and 12mm of sample] (easy) {Easy anchors\\\footnotesize levels 1--2};
\node[iodata, below right=8mm and 12mm of sample] (hard) {Hard anchors\\\footnotesize levels 3--4};

\node[stage, right=14mm of easy] (identity) {Identity\\\footnotesize no generation};
\node[block, right=14mm of hard] (gen) {DeepSeek\\\footnotesize Generator (DSPy)\\\footnotesize 5 candidates, batched};
\node[decision, right=18mm of gen] (val) {equivalent (DSPy) \&\\easy band (MARBERT)?\\\footnotesize best of 5 kept};
\node[block, above right=4mm and 10mm of val] (discard) {Discard};

\node[iodata, below right=14mm and 8mm of identity] (merge) {Combined\\training set};
\node[stage, right=14mm of merge] (finetune) {Fine-tune AraT5};

\draw[myarrow] (barec) -- (drop5);
\draw[myarrow] (drop5) -- (sample);
\draw[myarrow] (sample) -- (easy);
\draw[myarrow] (sample) -- (hard);
\draw[myarrow] (easy) -- (identity);
\draw[myarrow] (hard) -- (gen);
\draw[myarrow] (gen) -- (val);
\draw[myarrow] (val) -- node[lbl,pos=0.4,above,sloped]{no} (discard);
\draw[myarrow] (val.south) |- node[lbl, pos=0.7, above]{yes} (merge.east);
\draw[myarrow] (identity.south) |- (merge.north);
\draw[myarrow] (merge) -- (finetune);
\end{tikzpicture}}
\caption{Stage~1 training-data pipeline (current design). Level~5 is dropped
before sampling; a fair, stratified sample splits into easy anchors (identity
pairs, zero generation cost) and hard anchors (the only sentences that reach
the DeepSeek generator and the equivalence/level checks described in
Section~\ref{sec:qc}).}
\label{fig:datapipeline}
\end{figure}

\proscons{
  \item Roughly half the sampled anchors (the easy band) need zero LLM calls
  --- identity pairs are free and risk-free by construction, cutting both
  generation cost and the surface area for generation-introduced errors
  roughly in half versus generating for every sampled sentence.
  \item Dropping level~5 removes the riskiest single transformation
  (a direct expert-to-easy rewrite) while costing only 6.9\% of the corpus's
  \texttt{Specialized}-tagged content --- a favorable, data-verified trade,
  not a guess.
  \item Targeting a band rather than a specific sub-level avoids asking the
  generator (and the validator) to resolve a distinction that BAREC's own
  annotation scheme treats as fuzzy in the first place.
  \item BAREC is fully open (CC~BY-SA~4.0), confirmed directly (not gated,
  no agreement required) --- removing the license risk from the critical
  path (Section~\ref{sec:why-barec}).
  \item The easy-band identity pool needs no filtering or other
  preprocessing before use --- fragments, headers, and citations are real
  content a reader encounters, and identity pairs are safe by construction
  regardless of what the input looks like, so the pool is used exactly as
  BAREC provides it.
}{
  \item Every hard-band training pair is still LLM-generated; there is no
  ground-truth human-written rewrite anywhere in the pipeline, which keeps
  the full weight on the equivalence validator and level classifier
  working well.
  \item Both remaining DSPy programs (generator, equivalence validator)
  currently share one model (DeepSeek) for simplicity --- a model can share
  blind spots with itself across roles, so a generation error and the
  equivalence check that should catch it are not fully independent. The
  level check no longer carries this risk: it moved off DeepSeek entirely
  onto a separately fine-tuned local classifier (below). Worth periodically
  cross-checking a sample against a second model for the two that remain.
  \item Only a 2\% pilot sample (1,350 sentences: 708 easy / 606 hard / 36
  scripture-sourced sentences routed to identity pairs regardless of level,
  Section~\ref{sec:scripture}) has fully completed so far, with a 10\%
  sample now running against the reranked design --- real
  generation/validation pass-rates on the full hard-band pool (31,744
  sentences) are still not yet known.
}

\subsubsection{Why BAREC instead of SAMER}
\label{sec:why-barec}

The original version of this plan used the \textbf{SAMER} corpus
(arXiv:2404.18615) as the primary training-pair source, because it offers real,
human-written parallel rewrites. Checking SAMER's actual terms of use during
this revision surfaced a real problem: its license requires signing an NYU Abu
Dhabi usage agreement, restricts use to internal research/evaluation only,
prohibits redistribution or sublicensing, and requires a separate commercial
license for anything beyond that\footnote{Per the corpus's official project
page: \url{https://camel.abudhabi.nyu.edu/samer-simplification-corpus/}.} ---
real friction for a program where submitted work may need to be shared or
reused. SAMER is also narrow in domain: 159K words drawn entirely from 15
Arabic fiction novels, mostly published between 1865 and 1955.

\textbf{BAREC} (arXiv:2502.13520, ACL 2025 Findings) resolves the license
problem outright: it is released under CC~BY-SA~4.0, openly downloadable, no
agreement required.\footnote{Structure and license confirmed via the paper:
\url{https://arxiv.org/html/2502.13520v2}.} It is also an order of magnitude
larger (69,441 sentences, 1M+ words, 1,922 source documents) and far more
domain-diverse (education, literature, media, religion, reference works),
graded across 19 fine-grained readability levels rather than SAMER's two. The
catch: BAREC is \emph{not} a parallel corpus. Each sentence is independently
sourced and carries a single readability-grade label --- there are no aligned
``same content, two difficulty levels'' rewrites the way SAMER provides, so it
cannot directly substitute as ready-made (complex, simple) training pairs.

The resolution: use high-level BAREC sentences as the real, naturally-occurring
``complex'' side of every training pair, have Gemini generate the ``simple''
rewrite, and let the DSPy judge (Section~\ref{sec:qc}) filter every generated
pair on \emph{both} semantic equivalence and a measurable BAREC-calibrated
readability-level drop --- so a rewrite that merely rephrases without actually
simplifying gets rejected too.

\proscons{
  \item Fully open license --- no signed agreement, no redistribution
  restriction, no commercial-use gate. Removes a real legal/logistical risk
  from the critical path.
  \item Larger and far more domain-diverse than SAMER, so the ``complex'' side
  of every generated pair can be genuine, naturally-occurring Arabic from many
  registers, not confined to century-old literary fiction.
  \item 19 fine-grained levels provide a calibrated readability signal, usable
  as a second automatic filter beyond semantic equivalence alone.
}{
  \item Not a parallel corpus --- provides no aligned rewrites, so it cannot
  supply ready-made (complex, simple) training pairs the way SAMER did.
  \item Loses SAMER's free positive examples for the DSPy judge's validation
  set --- resolved by annotating the pilot's own real generated pairs
  instead of writing hand-labeled examples from scratch
  (Section~\ref{sec:annotation}).
  \item Some source genres (religious text, formal reference material) may be
  a poor stylistic match for a conversational simplification target and may
  need filtering by genre before sampling.
}

\subsubsection{Choosing the generation and judge models}
\label{sec:model-choice}

Two different model roles need two different selection criteria: the
\textbf{generation} model (turns BAREC sentences into candidate simplifications,
one bulk pass) is chosen for Arabic quality, since cost turns out to be
negligible regardless of vendor; the \textbf{judge} model (scores every
generated pair, plus the whole eval set --- much higher call volume) is where
cheap or self-hostable open models genuinely earn their place.

\paragraph{Generation model: Gemini.} Gemini's Arabic output quality has a
strong reputation, and at this data volume every vendor is cheap in absolute
terms, so quality --- not price --- should drive the choice. Estimating cost
for processing the \emph{entire} BAREC corpus ($\sim$1M+ words, $\sim$2M input
tokens and $\sim$2M output tokens at a working estimate of $\sim$2
tokens/Arabic word):

\begin{center}
\small
\begin{tabular}{@{}lrrr@{}}
\toprule
\textbf{Model} & \textbf{Input \$/M tok} & \textbf{Output \$/M tok} & \textbf{Est.\ cost, full corpus} \\
\midrule
Gemini 2.5 Flash-Lite\textsuperscript{\dag} & \$0.10 & \$0.40 & $\sim$\$1.00 \\
Gemini 3.5 Flash-Lite & \$0.30 & \$2.50 & $\sim$\$5.60 \\
Gemini 3.8 Flash (current default) & \$0.75 & \$3.75 & $\sim$\$9.00 \\
Gemini 3.1 Pro (highest Arabic quality) & \$2.00 & \$12.00 & $\sim$\$28.00 \\
\midrule
GPT-5.6 Luna (cheapest OpenAI tier) & \$0.20 & \$1.20 & $\sim$\$2.80 \\
GPT-5.6 Terra & \$2.00 & \$12.00 & $\sim$\$28.00 \\
GPT-5.6 Sol & \$5.00 & \$30.00 & $\sim$\$70.00 \\
\bottomrule
\end{tabular}
\end{center}

{\footnotesize Prices as of September 2026;\footnote{Gemini pricing:
\url{https://developer.puter.com/tutorials/gemini-api-pricing/}. OpenAI
pricing: \url{https://www.cloudzero.com/blog/openai-pricing/}.} Google's Batch
API cuts all Gemini rows roughly in half. \textsuperscript{\dag}Google is
retiring Gemini~2.5~Flash-Lite on 16~October~2026 --- do not build around it
this late in the timeline.}

\textbf{Verdict: cost is not a real constraint.} Even the most expensive row
here costs \$70 to process the \emph{entire} corpus; a sprint only needs a
curated subset, so real spend is likely a few dollars. Two practical gotchas
regardless of vendor: batch many sentences per API call rather than one call
per sentence (per-call instruction overhead otherwise adds up), and disable
``thinking'' mode on reasoning-capable variants, since hidden reasoning tokens
bill as output and can silently inflate cost.

\begin{noteb}
\textbf{Decided for the pilot: DeepSeek for all three roles.} For the initial
2\% pilot (Section~\ref{sec:datapipeline}), the generator and the DSPy
equivalence validator (Section~\ref{sec:qc}) are implemented against a
single DeepSeek model rather than splitting generation and judging across
vendors; the level check runs locally
(fine-tuned MARBERT, Section~\ref{sec:qc}) and makes no API call at all.
Reasoning: at pilot scale ($\sim$1,350 sentences at 2\%, now a 10\% sample
in progress) --- 2 batched DeepSeek calls per hard-band sentence (one
generating 5 candidates, one scoring all 5 for equivalence), easy-band
identity pairs need none --- the cost gap between vendors is immaterial,
and one vendor
integration is simpler to stand up and debug than two for a first pass.
Worth revisiting once the pipeline is validated at full scale (31,744
hard-band sentences) --- see the con already logged in
Section~\ref{sec:datapipeline} about one model sharing blind spots across
the two roles it still fills.
\end{noteb}

\subsection{Stage 2: Diacritization (Tashkeel Restoration)}

\begin{figure}[h!]
\centering
\begin{tikzpicture}[node distance=10mm and 9mm]
\node[iodata] (in) {Simplified sentence\\\footnotesize no diacritics};
\node[stage, right=18mm of in] (model) {Diacritization model\\\footnotesize e.g.\ CATT, fine-tuned\\\footnotesize on Tashkeela};
\node[iodata, right=18mm of model] (out) {Fully vocalized sentence\\\footnotesize every letter marked};

\draw[myarrow] (in) -- (model);
\draw[myarrow] (model) -- (out);
\end{tikzpicture}
\caption{Stage 2 in isolation. Tashkeela was already scoped by the team for a
separate, shelved idea --- directly reusable here.}
\label{fig:stage2}
\end{figure}

\proscons{
  \item Tashkeela is a large, well-established diacritization corpus; the team
  already has prior familiarity with it from a shelved idea, so onboarding cost
  is close to zero.
  \item Diacritization is a well-studied Arabic NLP task with published
  baselines to compare against, which de-risks the ``does this even work''
  question relative to a from-scratch task.
  \item Feeding \emph{simplified} rather than original text into this stage
  reduces the syntactic complexity the diacritizer has to handle, likely
  improving accuracy versus diacritizing the raw input directly.
}{
  \item Diacritization errors are read out loud verbatim by Stage~3's TTS --- an
  incorrect vowel mark becomes an incorrect \emph{sound}, which is precisely the
  failure mode the whole product exists to prevent. Errors here are more
  user-visible (audible) than in most diacritization use cases.
  \item Off-the-shelf diacritizers are usually trained on Modern Standard Arabic
  news/religious text (Tashkeela itself leans classical/religious); domain shift
  to simplified, contemporary, possibly conversational sentences could hurt
  accuracy and would need validation.
  \item Adds a second model to maintain, quantize, and keep in sync with Stage~1's
  output distribution.
}

\subsubsection{Architecture decision: byte-level generative vs.\ character-level tagging}
\label{sec:diacritization-arch}

Plugging a generic word/subword-level seq2seq model --- AraT5 included --- into
this stage is a poor fit, for a specific and instructive reason. AraT5's
tokenizer vocabulary was learned from a training corpus that is almost
entirely \emph{undiacritized}, since that is what nearly all Arabic text on
the internet looks like. A word like \textarabic{مدرسة} is probably one clean
token in that vocabulary; the same word with tashkeel,
\textarabic{مَدْرَسَةٌ}, is a different character string that the tokenizer's
merge rules were never optimized for, so it typically gets shattered into
several small, rare subword fragments instead of one clean token. The model
then has to learn a messy mapping from ``one token in'' to ``several oddly
fragmented tokens out'' --- a much harder learning problem than the task
actually requires.

\textbf{Byte-level tokenization sidesteps this entirely.} A model like ByT5
(Google) uses no learned vocabulary at all --- its ``tokens'' are raw UTF-8
bytes, 256 possible values, no OOV, ever. This turns out to fit Arabic
diacritization unusually well: tashkeel marks are not modifications of a
letter but separate Unicode combining-mark codepoints inserted immediately
after the base letter, and both Arabic base letters and diacritic marks
encode as 2-byte UTF-8 sequences in the same Unicode block. At the byte level,
diacritization looks close to \emph{``here is a byte stream with certain
2-byte chunks removed --- put them back in the right spots''} --- a
near-identity transformation, rather than the freer rewriting a word-level
model faces. This is consistent with Fine-Tashkeel's reported result: SOTA
accuracy from fine-tuning ByT5 with minimal training and no feature
engineering.

\textbf{The cost.} Byte sequences run roughly 2$\times$ longer than the
equivalent subword-tokenized sequence for Arabic, and since seq2seq decoding
is autoregressive, longer sequences mean slower generation. ByT5's own
architecture addresses this directly by deliberately shifting capacity toward
a heavier encoder and a lighter decoder, since decoder steps are where the
sequence-length cost bites hardest --- a real design trade-off worth factoring
into Stage~2's latency budget given the on-device framing (Section~\ref{sec:deployment-note}).

\begin{noteb}
\textbf{A nuance that cuts the other way from Section~\ref{sec:why-arat5-decision}.}
Even at the byte level, a seq2seq model is still \emph{generating} its output
--- nothing structurally prevents it from occasionally emitting a different
base letter than the input had; it is just empirically rare because the
transformation is so close to identity. A character-level \emph{tagging}
model (classify each existing letter into one of $\sim$15 diacritic classes,
as CAMeLBERT-style diacritizers do) cannot touch the base letters at all, by
construction --- it only ever attaches a label to a position that already
exists. That is the opposite lesson from the AraT5-vs-decoder-only debate for
Stage~1: there, the more constrained transduction architecture was the safer,
more faithful choice; here, a discriminative tagger is structurally the safer
choice over a generative seq2seq model, for the same underlying reason.

And yet, empirically, generative seems to win anyway: CATT's own published
results show its encoder-decoder (generative) checkpoint outscoring its
encoder-only (taggable) checkpoint on accuracy --- likely because an
autoregressive decoder can model dependencies between adjacent diacritic
choices that a plain per-position classifier cannot see the way a sequential
decoder does. \textbf{Practical takeaway:} the generative checkpoint is the
right default for accuracy, but because its faithfulness is empirical rather
than structural, it specifically warrants the base-letter-preservation check
called out as risk~R6 (Table~\ref{tab:risks}) once fine-tuned on Stage~1's
simplified-text domain --- that check is not automatically free the way it
would be with a tagger.
\end{noteb}

\subsubsection{Candidate models: CATT and Fine-Tashkeel}

\begin{center}
\small
\begin{tabular}{@{}p{2.5cm}p{4.75cm}p{4.75cm}@{}}
\toprule
& \textbf{CATT} (arXiv:2407.03236, ACL ArabicNLP 2024) & \textbf{Fine-Tashkeel / ByT5} (arXiv:2303.14588) \\
\midrule
Architecture & Character-level; ships \emph{both} an encoder-decoder (ED) and
an encoder-only (EO) checkpoint & Byte-level seq2seq; fine-tunes Google's
pretrained ByT5 \\
\addlinespace
License & Apache 2.0 --- relicensed from CC-BY-NC specifically to enable
broader adoption & Not directly confirmed here --- check the paper's linked
repo before relying on it \\
\addlinespace
Reported result & SOTA on their benchmark; beats GPT-4-turbo by a relative
9.36\% diacritic error rate & SOTA claim; $\sim$40\% WER reduction vs.\
baselines, minimal training, no feature engineering \\
\addlinespace
Ready to fine-tune & Yes --- pip-installable (\texttt{catt-tashkeel}), ships
\texttt{train\_catt.py} & A recipe more than a shipped checkpoint --- would
mean fine-tuning a public ByT5 checkpoint yourself following the paper's
method \\
\addlinespace
Best-fit note & Concrete, ready starting point: use the ED checkpoint for
accuracy, EO for speed if on-device latency is tight & Fallback approach if
CATT underperforms after domain adaptation to simplified text \\
\bottomrule
\end{tabular}
\end{center}

\textbf{Recommendation:} start from CATT's encoder-decoder checkpoint ---
concrete weights, a permissive license, and a ready fine-tuning script beat a
recipe that has to be reproduced from a paper. Fine-tune it further on
Tashkeela plus a sample of Stage~1's simplified output specifically (not just
Tashkeela's native classical/religious domain, per risk~R6), and validate
base-letter preservation explicitly given the faithfulness nuance above,
before treating its output as trustworthy enough to read aloud verbatim in
Stage~3. A recent shared task, KSAA-2026 (OSACT7~@~LREC~2026), ran a
systematic 18-model bake-off across seq2seq, token-classification,
decoder-LLM, and ASR-based approaches for this exact task --- worth a glance
at the leaderboard for the freshest signal before committing.

\begin{noteb}
\textbf{Aside, outside this plan's scope: a tokenizer-free tagging
alternative.} CANINE (Clark et al., TACL 2021) is a genuine tokenization-free
BERT-style encoder --- no vocabulary at all, it hashes raw Unicode codepoints
directly --- and a natural fit for Arabic's non-concatenative
(root-and-pattern) morphology, where there is no clean contiguous split point
for a subword tokenizer to find. Paired with a per-character classification
head, it would give a tokenizer-free CAMeLBERT analog that structurally
cannot alter base letters (Section~\ref{sec:diacritization-arch}) --- CATT's
own encoder-only checkpoint is already effectively this idea in practice,
built from an unnamed ``pretrained character-based BERT.'' A lower-effort
variant: take ByT5's pretrained encoder alone, swap the decoder for a
classification head, and fine-tune as a tagger. Neither has a confirmed
Arabic-diacritization benchmark and neither is part of this capstone's
recommendation --- noted here purely as a standalone idea worth exploring
separately.
\end{noteb}

\subsection{Stage 3: Read Aloud with Synced Highlighting}
\label{sec:tts}

\begin{noteb}
\textbf{This stage is optional.} Unlike Stages~1--2, Stage~3 is not required
to ship a defensible capstone: a text-only product (typed input
$\to$ simplified $\to$ diacritized output, no audio) is already a complete,
demoable accessibility tool on its own. Section~\ref{sec:tts-arch} below lays
out a genuine three-way trade-off between voice quality, license, and native
word-timing support, with \emph{no clean winner} among the options found ---
if that trade-off resolves badly (e.g.\ no adequately open, commercially
usable option is found in time), the right call for a fast 3-week sprint is
to cut Stage~3 entirely rather than force a compromised implementation,
freeing that time for Stages~1--2 and the QC pipeline instead.
\end{noteb}

\begin{figure}[h!]
\centering
\resizebox{\textwidth}{!}{%
\begin{tikzpicture}[node distance=10mm and 9mm]
\node[iodata] (in) {Vocalized sentence};
\node[stage, right=18mm of in] (tts) {Arabic TTS engine};
\node[iodata, above right=6mm and 16mm of tts] (audio) {Audio waveform};
\node[iodata, below right=6mm and 16mm of tts] (timing) {Word/phoneme\\timestamps};
\node[block, right=32mm of tts] (sync) {Highlight sync engine};
\node[iodata, right=18mm of sync] (ui) {UI: karaoke-style\\word highlighting\\+ audio playback};

\draw[myarrow] (in) -- (tts);
\draw[myarrow] (tts) -- (audio);
\draw[myarrow] (tts) -- (timing);
\draw[myarrow] (audio) -| (sync);
\draw[myarrow] (timing) -| (sync);
\draw[myarrow] (sync) -- (ui);
\end{tikzpicture}}
\caption{Stage 3 in isolation. The critical dependency is a TTS engine that
exposes (or can be made to expose) word-level timestamps, not just raw audio.}
\label{fig:stage3}
\end{figure}

Word-by-word highlighting synced to audio (``karaoke-style'' reading support) is a
well-established multi-sensory reading technique independent of language, so the
main engineering question is purely which Arabic TTS engine to use, and whether it
gives you timing information directly or you have to derive it via forced alignment.

\begin{center}
\small
\begin{tabular}{@{}p{3.2cm}p{3.5cm}p{3.5cm}p{3.5cm}@{}}
\toprule
& \textbf{Cloud (Azure / Google / Amazon Polly Arabic Neural)} &
\textbf{Open-source (e.g.\ Coqui/VITS Arabic, Meta MMS-TTS)} &
\textbf{On-device OS TTS (Android/iOS Arabic voices)} \\
\midrule
Voice quality & Highest, near-natural & Variable, improving fast & Decent, OS-dependent \\
Word timestamps & Often exposed natively (SSML marks) & Depends entirely on architecture --- see Section~\ref{sec:tts-arch} & Inconsistent API support \\
On-device / offline & No (network required) & Yes, if quantized & Yes, built-in \\
Cost & Per-character billing & Free (compute cost only) & Free \\
Setup effort (3-week sprint) & Low & Medium--High & Low--Medium \\
\bottomrule
\end{tabular}
\end{center}

\proscons{
  \item A cloud neural TTS (Azure/Google/Polly all have solid Arabic voices) gets
  a working, good-quality demo running fastest --- important given the 3-week
  clock.
  \item SSML-based engines often return word or viseme timing directly, avoiding
  a whole forced-alignment sub-problem.
  \item This stage is largely orthogonal to the NLP work in Stages~1--2, so it
  can be built and demoed in parallel by a different team member.
}{
  \item Best cloud Arabic voices conflict with the program's on-device framing
  unless deliberately scoped as ``Stage~3 is a non-on-device convenience layer,
  the on-device claim rests on Stages~1--2.''
  \item Open-source/on-device Arabic TTS with reliable word-level timing is
  genuinely less mature than the Latin-script equivalent --- this is a real,
  not-yet-fully-solved sub-problem, not just an integration task.
  \item If timestamps must be derived via forced alignment rather than read
  directly from the engine, that is additional, nontrivial engineering scope
  layered on top of an already three-stage pipeline.
}

\subsubsection{Architecture decision: native alignment vs.\ bolted-on forced alignment}
\label{sec:tts-arch}

The same analytical move that shaped Stages~1 and~2 applies here: word-level
timing is either a native byproduct of a TTS model's own generation
mechanism, or it is something bolted on afterward with a separate
forced-alignment pass --- and which category a given model falls into is a
direct consequence of its architecture, not an accident of implementation
quality.

\begin{itemize}[leftmargin=1.4em]
  \item \textbf{VITS-family models (native alignment).} VITS is trained with
  an internal monotonic alignment search that learns a genuine
  text-to-audio-frame duration mapping during training; at inference, a
  duration predictor outputs how many frames each phoneme occupies
  \emph{before} the decoder generates any audio. Per-phoneme --- and by
  extension per-word --- start/end times fall out of generation directly, no
  extra step required. Meta's MMS-TTS and Piper are both VITS-based.
  \item \textbf{FastSpeech2 / Matcha-TTS / StyleTTS2 family (native at
  inference, cost paid at training time).} Also has an explicit duration
  predictor, giving the same inference-time benefit --- but this family
  specifically needs externally-supplied ground-truth durations (from a
  forced aligner such as the Montreal Forced Aligner, or from a pretrained
  autoregressive teacher model) \emph{during training} to learn what to
  predict. That cost is paid once, up front, not per inference call. Nabra-82M
  (Section~\ref{sec:tts-models}), a StyleTTS2 fine-tune, is a concrete
  example in this bucket.
  \item \textbf{Flow-matching without a duration model (F5-TTS family, and
  by extension SILMA TTS, which is built on it).} F5-TTS deliberately drops
  the explicit duration predictor to simplify its pipeline, estimating total
  utterance duration from a crude length-ratio heuristic (audio-prompt
  transcript length vs.\ target text length) rather than predicting a
  duration per phoneme or word. There is no reliable per-word timing signal
  to read off the model at all --- recovering one requires exactly the
  bolted-on forced-alignment step this whole analysis is trying to avoid.
  \item \textbf{Purely autoregressive attention-based synthesis} (classic
  Tacotron-style, or any token-based/LLM-style generative TTS with no
  explicit duration model). Timing information, if present at all, is
  buried in attention weights that are not reliable position markers ---
  attention can skip, repeat, or drift. The worst case for this pipeline's
  needs.
\end{itemize}

\begin{noteb}
\textbf{The tension this creates.} The architecture family that is
structurally correct for the timing requirement (VITS) is not the family
currently producing the highest-quality Arabic voices (F5-TTS-family models
like SILMA TTS report state-of-the-art naturalness). Unlike Stages~1--2,
where the structurally-correct choice and the highest-quality choice mostly
lined up, Stage~3 genuinely pulls in two directions at once --- see the
candidate comparison below.
\end{noteb}

\subsubsection{Candidate models}
\label{sec:tts-models}

\begin{center}
\small
\begin{longtable}{@{}p{2.3cm}p{1.7cm}p{1.5cm}p{2.3cm}p{1.6cm}p{3.7cm}@{}}
\toprule
\textbf{Model} & \textbf{Arch.} & \textbf{Params} & \textbf{License} &
\textbf{Native timing?} & \textbf{Notes} \\
\midrule
\endhead
SILMA TTS~v1 & F5-TTS (flow matching, DiT) & 150M & Apache~2.0 & No ---
length-ratio heuristic only & Best-in-class Arabic-specific quality; RTF~$\sim$0.12
on an RTX~4090 \\
\addlinespace
Nabra-82M (oddadmix) & StyleTTS2 (Kokoro-82M fine-tune) & 82M & Apache~2.0 &
Yes, unconfirmed --- has a duration predictor, but not confirmed exposed as
an output in the released demo/code & Dedicated phoneme slots for
ع/ح that generic phonemizers collapse; \emph{expects diacritized input} ---
a direct fit with Stage~2's output, no redundant preprocessing; MLX build
for fully offline Apple Silicon/iOS; RTF~$<$1 on CPU; single voice
(\texttt{af\_msa}, MSA female) only; built and maintained by one independent
developer, not a funded lab --- newer and less battle-tested than the other
rows here \\
\addlinespace
MMS-TTS-Ara (Meta) & VITS & 36.3M & CC-BY-NC-4.0 (non-commercial) & Yes ---
native duration predictor & Smallest option; quantized ONNX build already
exists for browser/edge (\texttt{Xenova/mms-tts-ara}); one of $\sim$1{,}100
MMS languages, not Arabic-specialized \\
\addlinespace
Piper (``kareem'', ar\_JO) & VITS & Tens of millions (not confirmed exactly)
& Piper itself MIT; per-voice license varies & Yes --- native duration
predictor & Most edge-proven runtime (real-time on a Raspberry~Pi~4 across
100+ voices); Arabic coverage currently thin (one low/medium-quality voice
plus two unofficial community voices); ships its own tashkeel step ---
disable it, since Stage~2 already supplies diacritized text \\
\addlinespace
ArTST (MBZUAI) & SpeechT5 & Not confirmed here & CC-BY-NC-4.0
(TTS checkpoint) & Not confirmed here & Pretrained from scratch on Arabic
MSA speech/text, not a multilingual model adapted to Arabic; same
non-commercial-license pattern already seen with SAMER (Section~\ref{sec:why-barec}) \\
\addlinespace
Cloud (Azure / Google / Polly) & Proprietary & N/A & Per-vendor commercial
terms & Yes --- SSML word-boundary / viseme events & Highest current
naturalness; not on-device; per-character billing \\
\bottomrule
\end{longtable}
\end{center}

\textbf{Nabra-82M meaningfully improves this trade-off, without fully
resolving it.} It is the only option here that combines a fully permissive
license (Apache~2.0), an architecture family that natively supports duration
prediction (StyleTTS2, not flow-matching-without-duration like SILMA), a
genuine offline/edge story (MLX build for Apple Silicon/iOS), and --- uniquely
--- an input contract that matches this pipeline exactly: it \emph{expects}
diacritized text, which is precisely what Stage~2 already produces, with no
redundant or conflicting preprocessing step to disable (unlike Piper's
built-in tashkeel step).

\begin{itemize}[leftmargin=1.4em]
  \item \textbf{Nabra-82M} is the strongest combination on paper, but carries
  real newness risk: whether word-level timing is actually exposed (versus
  just architecturally present) needs a hands-on check before relying on it,
  it ships one voice only, and it is maintained by a single independent
  developer rather than a funded lab or large open-source project ---
  meaningfully less battle-tested than SILMA, MMS-TTS-Ara, or Piper. Pilot it
  first, but do not commit to it sight-unseen.
  \item \textbf{MMS-TTS-Ara} remains the most-established ``cheap and
  structurally correct'' fallback: smallest footprint, an existing quantized
  edge build, and confirmed native timing --- but its CC-BY-NC-4.0 license
  blocks commercial use without separate permission from Meta, echoing the
  exact SAMER lesson from Section~\ref{sec:why-barec}: check the license
  before building around a model, not after.
  \item \textbf{Piper} is the most genuinely edge-proven runtime, but
  Arabic voice quality/coverage is currently the weakest link among the
  options considered.
  \item \textbf{SILMA TTS~v1} has the best-documented Arabic-specific voice
  quality and a fully commercial-friendly license, but its F5-TTS lineage
  means word timing is not a native output --- using it means reintroducing
  the bolted-on forced-alignment step this analysis exists to avoid.
\end{itemize}

This is a genuine team decision (Section~\ref{sec:open-items}), not a
default: pilot Nabra-82M first given how well it fits on paper, fall back to
MMS-TTS-Ara or Piper if its timing turns out not to be usable or its
newness proves too risky for the sprint timeline, and fall back further to
SILMA TTS plus a forced-alignment step if voice quality and
commercial-friendliness matter more than native timing.

\begin{noteb}
\textbf{This is still an actively explorable field --- treat this table as a
snapshot, not a closed decision.} Arabic TTS is moving fast enough that a
genuinely strong candidate (Nabra-82M) turned up mid-research and wasn't on
the radar when this section was first drafted; the community-run Arabic TTS
Arena leaderboard alone tracks 15+ systems being actively compared by native
listeners, and new entrants (several tied to the current wave of Saudi
voice-AI startups) keep appearing. No model evaluated here is a confirmed
100\% fit --- each carries at least one open question (unexposed timing,
non-commercial license, thin voice coverage, or unproven maturity). Re-check
the landscape close to when Stage~3 is actually built, rather than locking
in this list today.
\end{noteb}

\newpage
% =====================================================================
\section{Quality-Control Pipeline: LLM-as-Judge + DSPy}
\label{sec:qc}
% =====================================================================

The pipeline defines a \textbf{generator} that produces 5 candidate
simplifications per hard-band anchor in one batched call, plus two
independent checks that gate them: a DSPy \textbf{equivalence validator}
and a locally hosted, fine-tuned \textbf{level classifier} --- one for
semantic equivalence, one for whether a candidate actually lands in the
easy band (Section~\ref{sec:datapipeline}). Splitting the level check out
from the meaning check keeps each program's job narrow and separately
debuggable: a single combined judge has to get both dimensions right in one
verdict, and it is harder to tell which dimension failed when a pair gets
rejected. The generator and equivalence validator run against DeepSeek; the
level classifier is not an LLM call at all (below). This is a change from
an earlier revision, where the level check was itself a third DSPy program
(\texttt{ClassifyReadabilityLevel}) prompted against DeepSeek --- superseded
once BAREC's $\sim$69k gold labels made supervised fine-tuning both cheaper
and more accurate than few-shot prompting (Section~\ref{sec:next-phase}).

\begin{noteb}
\textbf{Revised since the last pass: the level check left DSPy for a
fine-tuned local classifier; the equivalence validator's signature is now
batched for reranking.} The generator (\texttt{SimplifyToEasyArabic}) still
has no ground truth to compile against --- no human-written simplification
exists anywhere in BAREC or elsewhere in this pipeline --- so it still runs
zero-shot by necessity, now producing 5 candidates per call instead of 1
(Section~\ref{sec:next-phase}).

\textbf{Level classification: fine-tuned MARBERT, quantized to ONNX
int8} (\texttt{scripts/train\_level\_classifier\_bert.py}, weights at
\texttt{models/level\_classifier\_bert\_marbert\_onnx\_int8/}), replacing
the DSPy-prompted \texttt{ClassifyReadabilityLevel} entirely. BAREC's
$\sim$69k gold level labels are enough to supervise-train a classification
head directly, at far greater scale than the $\sim$1,000-sentence pool the
old few-shot classifier was compiled against, with zero manual labeling and
no per-call API spend at inference. Several architectures were benchmarked
against a real held-out split before settling on this one: a 4-way softmax
head on CAMeLBERT-mix scored 83.5--83.9\% band accuracy; binary,
CORAL-ordinal, and hybrid heads on the same base landed at 83.4\%, 80.1\%,
and 83.0\% respectively --- architecture changes on one encoder converged
to roughly the same ceiling. Swapping the base model instead of the head
is what actually moved the number: AraBERT matched CAMeLBERT-mix at
83.4\%, but \textbf{MARBERT reached 86.8--87.5\% band accuracy}, 3--4pp
above the entire CAMeLBERT-mix-based cluster, attributed to its
pretraining corpus (Twitter-scale dialectal/MSA Arabic) rather than any
architectural difference. A majority-vote ensemble of CAMeLBERT-mix +
AraBERT + MARBERT was also tried, on the theory that three models would
catch each other's mistakes; it scored 84.8\%, \emph{worse} than MARBERT
alone --- diagnosed as 31 net-negative overrules against only 15
net-positive ones, i.e.\ the weaker two models outvoted the better one more
often than they corrected it. MARBERT alone is what ships, quantized
dynamic-int8 via ONNX (164MB, CPU-only inference through
\texttt{onnxruntime} directly --- no \texttt{torch} dependency at inference
time, since the tensor glue an \texttt{optimum} wrapper would add is
unnecessary when the ONNX graph and tokenizer output can run as plain
NumPy arrays).

\textbf{\texttt{CheckSemanticEquivalence}} (\texttt{scripts/optimize\_equivalence\_validator.py}):
still the one validator with no BAREC-native ground truth --- the
annotation workflow described below is what supplies it. The signature
itself changed to support reranking: it now takes a \emph{list} of
candidate simplifications and returns a list of equivalence scores plus a
short reasoning string per candidate (previously discarded when only the
overall DSPy trace was inspected; now a first-class output field, both
because it is directly useful to a human reviewer and because
\texttt{dspy.Predict} with its own explicit reasoning field made
\texttt{dspy.ChainOfThought}'s implicit one redundant). Compiled with
\texttt{dspy.MIPROv2} against the same 102 human-annotated pairs (65
equivalent / 37 not-equivalent), split 77 train / 25 held-out, adapted to
the new batched shape as batch-of-1 calls (one candidate per list) since no
genuine multi-candidate ground truth exists yet --- a stated limitation of
this evaluation, not of production use, which always batches 5 real
candidates together. \textbf{Result: mean score closeness improved 0.652
$\to$ 0.668, but thresholded accuracy (score $\geq$0.5 as ``equivalent'')
was flat at 64.0\% both before and after compiling.} That number is less
reassuring than it looks: at threshold 0.5, both the raw and the
MIPROv2-compiled validator scored \emph{every one} of the 25 held-out
pairs $\geq$0.5 (16/16 true positives, but also 9/9 false positives, 0
true negatives) --- 64.0\% is exactly the positive-class base rate
(16/25), not evidence of a working negative classifier, a real regression
in kind (not just degree) from the earlier single-candidate signature,
which correctly rejected 7--8 of 9 negatives. Whether this is a genuine
behavior change or an artifact of scoring one candidate in isolation
inside a signature designed for five (Section~\ref{sec:next-phase}) is not
yet known and is flagged here rather than smoothed over; re-evaluating
once the 10\% pilot's real 5-candidate batches accumulate enough annotated
negatives is the next step. Separately, \textbf{a cheaper alternative was
tested and rejected}: BERTScore F1 between original and candidate showed
essentially zero correlation with the human equivalence labels
($r=-0.014$) and a calibrated decision threshold reached only 48--64\%
accuracy, no better than (and sometimes worse than) the majority-class
baseline --- confirming an LLM judge is doing real work here, not
something a cheap embedding-similarity metric can substitute for.
Enabling DeepSeek's ``thinking'' mode on the judge was also tried and made
results worse (0.628 vs.\ 0.651 mean closeness on an earlier evaluation
pass, including one case where hidden reasoning consumed the entire
4096-token budget before any visible answer was written) --- thinking
stays disabled, consistent with Section~\ref{sec:deepseek-lessons}.
\end{noteb}

\begin{figure}[h!]
\centering
\resizebox{\textwidth}{!}{%
\begin{tikzpicture}[node distance=9mm and 9mm]
\node[iodata] (hard) {Hard anchor\\\footnotesize \texttt{original\_text}};
\node[stage, right=16mm of hard] (gen) {\texttt{Generator}\\\footnotesize $\to$ 5 $\times$ \texttt{simplified\_text}};
\node[iodata, right=16mm of gen] (cand) {5 candidate pairs};

\node[stage, below right=10mm and 4mm of cand] (eq) {\texttt{CheckSemanticEquivalence}\\\footnotesize (original, [5 candidates])\\\footnotesize $\to$ \texttt{equivalence\_scores}, \texttt{reasoning}};
\node[stage, above right=10mm and 4mm of cand] (diff) {Local MARBERT classifier\\\footnotesize (each candidate)\\\footnotesize $\to$ \texttt{level: int} (1--4)};

\node[decision, right=30mm of cand] (gate) {best candidate:\\equivalent \&\\level $\leq$ 2?};
\node[iodata, right=16mm of gate] (kept) {Kept pair};
\node[block, above right=8mm and 4mm of gate] (rej) {Discard};

\draw[myarrow] (hard) -- (gen);
\draw[myarrow] (gen) -- (cand);
\draw[myarrow] (cand) -- (eq);
\draw[myarrow] (cand) -- (diff);
\draw[myarrow] (eq) -- (gate);
\draw[myarrow] (diff) -- (gate);
\draw[myarrow] (gate) -- node[lbl,above]{yes} (kept);
\draw[myarrow] (gate) -- node[lbl,pos=0.4,above,sloped]{no} (rej);
\end{tikzpicture}}
\caption{The generator, DSPy equivalence validator, and local MARBERT
classifier, and how the best of 5 reranked candidates is kept or
discarded. Both checks must pass independently for the winning candidate
--- a fluent but meaning-distorted rewrite is caught by
\texttt{CheckSemanticEquivalence} even if the level classifier reports an
easy level, and vice versa. \texttt{CheckSemanticEquivalence} now declares
an explicit \texttt{reasoning\_per\_candidate} output field alongside its
scores, not just a bare verdict; it is DSPy-optimizer-compiled against real
ground truth (results below), the level classifier is separately
fine-tuned and quantized (also below), and the generator remains zero-shot
for lack of any ground truth to compile against.}
\label{fig:dspy}
\end{figure}

\subsection{Where the pipeline is used}

\begin{enumerate}[leftmargin=1.6em]
  \item \textbf{Data filtering (build-time, hard band only).} Every hard-band
  anchor's generated candidate is scored by both validators before it enters
  the AraT5 training set; a pair is kept only if both pass
  (Figure~\ref{fig:datapipeline}). Easy-band identity pairs skip this
  entirely --- there is nothing to validate when the output is the input.
  \item \textbf{Evaluation (build-time).} The same two validators run on a
  held-out split; the average equivalence and easy-band pass rates are the
  headline faithfulness and effectiveness metrics for the pitch ---
  meaningfully stronger evidence than BLEU/ROUGE, which reward word overlap, not
  meaning preservation.
  \item \textbf{Runtime gate (optional, stretch goal).} Score live output before
  showing it to the user; if either validator fails, fall back to the original text
  with a note rather than silently showing a possibly-distorted simplification.
  This echoes a ``trust signaling'' pattern already used in a sibling team idea
  (the RAG curriculum assistant, which cites its sources).
\end{enumerate}

\proscons{
  \item Directly targets the metric that actually matters for this product
  (meaning preservation), rather than a proxy metric.
  \item DSPy's optimizer-compiled prompt is a defensible, reproducible
  methodology to present in a pitch --- ``we didn't hand-tune a prompt and hope,
  we optimized it against labeled data,'' which fits SIC's evident preference for
  rigor over vibes.
  \item Reusable across three pipeline touchpoints (filtering, eval, optional
  runtime gate) for one build cost.
  \item The runtime gate, if built, becomes a genuine user-facing safety feature,
  not just an internal QA tool.
}{
  \item Adds a second model-development track (the judge) on top of the main
  model (AraT5) inside the same 3-week sprint --- real scope, not free.
  \item DSPy optimization requires a hand-labeled validation set that does not
  yet exist; ``negative'' (meaning-distorted) examples specifically must be
  manually constructed, since real corpora are overwhelmingly positive examples.
  \item Judge reliability on Arabic semantic equivalence is empirically
  unverified going in --- this is a real open risk, not a formality (see below).
  \item A weak or miscalibrated judge silently degrades everything downstream of
  it: bad filtering lets bad training data through, and a bad eval score either
  falsely reassures the team or falsely tanks a model that was actually fine.
}

\subsection{Human annotation: how the A3 validation set actually got built}
\label{sec:annotation}

The original plan for A3 (Section~\ref{sec:open-items}) was to hand-construct
a labeled set from scratch: manually simplify some BAREC sentences for
positives, and deliberately corrupt others (drop a negation, swap a number,
substitute an antonym) for negatives, since BAREC provides no aligned
rewrites and therefore no free positive examples the way SAMER's
human-written simplifications would have. That plan is superseded by
something better: instead of synthetic corrupted examples, the pilot's own
\emph{real} generated pairs are annotated directly. A web tool (``Pilot Pair
Review,'' published as a Claude Artifact with a shared realtime database) lets
a reviewer browse every generated pair --- original and simplified text
side by side, with RTL Arabic typography --- and record two independent
verdicts per pair: \texttt{level\_verdict} (\emph{ok} / \emph{should\_be\_easy}
/ \emph{should\_be\_hard}) and \texttt{equiv\_verdict} (\emph{ok} /
\emph{should\_be\_equivalent} / \emph{should\_not\_be\_equivalent}), plus a
free-text note. This is real ground truth on real pipeline output, not
synthetic corruption of hand-picked examples, and it is what the equivalence
validator's 102-pair calibration set (Section~\ref{sec:qc}) is drawn from.
Annotation is ongoing; more pairs (especially clear level corrections, which
are thinner than the equivalence corrections so far) will continue to
sharpen both validators.

\subsection{Content safety: scripture must be preserved verbatim}
\label{sec:scripture}

BAREC draws sentences from $\sim$30 source categories, four of which are
scripture: Quran, Hadith, Old Testament, and New Testament. A hard-band
scripture sentence is, by definition, exactly the kind of sentence the
generator is supposed to rewrite --- and rewriting sacred or canonical text
is never acceptable, regardless of how genuinely difficult it reads. This
was not a hypothetical risk: an early completed pilot run, before this fix
existed, sent 28 scripture-sourced sentences through the generator,
including a Quranic ayah with its verse number altered and a hadith with
the \textarabic{ﷺ} honorific reworded. The reviewer caught this by hand.

\textbf{Fix.} At sampling time, sentences whose \texttt{Source} column is
Quran, Hadith, Old Testament, or New Testament are filtered out of the
easy/hard split entirely and routed to identity pairs (output~=~input) ---
the same free, safe-by-construction mechanism already used for easy-band
sentences, applied here regardless of the sentence's assessed difficulty.
Validated by dry run: a 2\% stratified sample now produces 1,350 sentences
as 708 easy / 606 hard / 36 scripture (preserved verbatim) --- 28 fewer
hard-band sentences than before the fix, matching exactly the 28 scripture
sentences the reviewer had found in the earlier completed pilot.

\begin{noteb}
\textbf{Known gap, not yet solved: this fix only cleans the training data,
not deployment.} The fix above works because BAREC sentences carry
\texttt{Source} metadata to filter on. The deployed product takes raw
typed or pasted sentences as input, with no such metadata --- a real end
user pasting an unfamiliar Quranic or Hadith sentence would have nothing
for the pipeline to route on, and it could still be rewritten. A scoped,
not-yet-built follow-up: a Quran-specific runtime exact/fuzzy-match lookup
guard is genuinely tractable, since the Quran is a small, fixed, fully
enumerable text ($\sim$6,236 verses) --- unlike Hadith or the Bible, which
are far larger and less centrally indexed, and realistically out of reach
for full coverage within a 3-week capstone. This is future work, explicitly
flagged as such, not claimed as solved.
\end{noteb}

\subsection{Pipeline resilience: checkpointing and timeouts}
\label{sec:resilience}

A real incident shaped this: a $\sim$2-hour pilot run hung indefinitely at
600/634 hard-band sentences after the development laptop slept, silently
killing the underlying network connection with nothing to catch it ---
everything already processed would have been lost, since nothing was
persisted until the run finished. Two fixes followed. First, every LM call
carries a 90-second timeout, so a stalled connection fails fast and is
caught by the per-example error handling instead of hanging the whole run.
Second, \texttt{run\_hard\_pipeline} in
\texttt{scripts/barec\_simplification\_pipeline.py} writes every successful
result to a JSON-lines checkpoint file immediately (append, flush, and
fsync after each one, not held in memory until the run completes); on
restart, already-completed sentence IDs are filtered out of the remaining
work automatically, and a checkpoint from a different sample configuration
is detected and not silently reused.

\subsection{Engineering lessons: DeepSeek V4.1 Flash integration}
\label{sec:deepseek-lessons}

Three integration issues, all confirmed against the real API, are worth
recording so they are not rediscovered later:

\begin{itemize}[leftmargin=1.4em]
  \item \textbf{Thinking mode must be disabled}
  (\texttt{thinking=\{"type": "disabled"\}}). Left on, DeepSeek's internal
  reasoning pass consumed the entire token budget on nearly every call
  before the model ever wrote the structured answer DSPy needs, causing
  near-total call failure --- not just added cost, as a generic
  ``reasoning tokens bill as output'' concern might suggest, but functional
  breakage. \texttt{dspy.ChainOfThought}'s own \texttt{reasoning} field
  already serves this purpose within the visible token budget.
  \item \textbf{DSPy's JSON-adapter fallback must be disabled}
  (\texttt{dspy.ChatAdapter(use\_json\_adapter\_fallback=False)}). DSPy's
  default behavior retries a parse failure via a JSON-mode adapter, which
  sends a \texttt{response\_format} parameter DeepSeek's endpoint rejects
  outright --- turning a single recoverable parse failure into a fatal
  crash.
  \item \textbf{\texttt{optuna} is a required but not automatically
  installed dependency} of \texttt{dspy.MIPROv2} (its internal trial
  search uses Optuna) --- install it explicitly
  (\texttt{uv add optuna}) before running any MIPROv2 compile.
\end{itemize}

\subsection{Implemented: candidate reranking and a fine-tuned classifier}
\label{sec:next-phase}

The original design gated generation with a binary pass/fail: one candidate
per hard sentence, kept only if both validators passed. This has been
replaced with a generate-then-rerank design, directly informed by the
actual TSAR 2025 Shared Task on Readability-Controlled Text Simplification
(EMNLP 2025) --- its winning system, EhiMeNLP (Miyata, Horiguchi, Kondo,
Fujiwara, and Kajiwara), used exactly this two-step strategy: generate
multiple simplification candidates via LLM prompting, then rerank by
readability-based filtering and semantic similarity to the original, with
no parallel training data required. What actually shipped:

\begin{enumerate}[leftmargin=1.6em]
  \item \textbf{Fine-tuned a pretrained Arabic transformer} (MARBERT, after
  benchmarking it against CAMeLBERT-mix and AraBERT across four
  classification-head designs) with a classification head on the full
  BAREC corpus, replacing the LLM-prompted \texttt{ClassifyReadabilityLevel}
  outright rather than merely supplementing it. BAREC's $\sim$69k
  gold-labeled sentences, versus the $\sim$1,000 used for the old few-shot
  calibration, made supervised fine-tuning both cheaper (no per-call API
  spend at inference) and markedly more accurate: 86.8--87.5\% band
  accuracy versus 79.7\% for the DSPy-prompted predecessor. Full
  architecture search and numbers in Section~\ref{sec:qc}.
  \item \textbf{Replaced binary gating with generate-N-then-rerank}: the
  generator now produces $N=5$ candidate simplifications per hard sentence
  in one batched call, the equivalence validator scores all 5 in a second
  batched call, the local classifier scores all 5 for free, and the
  best-scoring candidate that passes both checks is kept --- 2 DeepSeek
  calls per hard sentence total, down from 3 in the original one-candidate
  design, despite generating 5$\times$ the candidates.
  \item \textbf{Stretch goal, still not built: retrieval-augmented
  generation.} Use the fine-tuned classifier's embeddings to retrieve
  same-topic easy-band BAREC neighbors as few-shot context for the
  generator --- reusing BAREC's own labels as a retrieval index, again with
  no parallel data required.
\end{enumerate}

The GPU blocker noted in an earlier revision (development machine has no
CUDA GPU) was resolved by renting GPU time on vast.ai for the fine-tuning
runs, not by acquiring local hardware --- a real, repeatable path for the
rest of the sprint if further fine-tuning is needed, with its own gotchas
(2FA/SSH setup, a CUDA forward-compatibility library conflict on one fresh
host, transformers API drift against older tutorials) now documented for
reuse rather than rediscovered each time.

\newpage
% =====================================================================
\section{End-to-End System Architecture: What Runs Where}
% =====================================================================

Because the pitch leans on an ``on-device AI'' framing, it matters to be precise
about which components actually need to run on-device at inference time versus
which are strictly build-time, offline tooling.

\begin{figure}[h!]
\centering
\resizebox{\textwidth}{!}{%
\begin{tikzpicture}[node distance=7mm and 10mm]

\node[cloudbox, minimum width=6.6cm, minimum height=5.6cm] (cloud) {};
\node[anchor=north, font=\bfseries, text=violet!60!black] at (cloud.north) {Cloud / Offline --- Build-Time Only};
\node[block, below=10mm of cloud.north] (bigllm) {Gemini\\\footnotesize synthetic data generation};
\node[block, below=6mm of bigllm] (judgebuild) {DSPy judge\\\footnotesize filtering + eval};
\node[block, below=6mm of judgebuild] (trainloop) {AraT5 fine-tuning loop};

\node[devicebox, right=26mm of cloud, minimum width=6.6cm, minimum height=5.6cm] (device) {};
\node[anchor=north, font=\bfseries, text=teal!60!black] at (device.north) {On-Device --- Runtime};
\node[stage, below=10mm of device.north] (aradev) {Quantized AraT5\\\footnotesize simplify};
\node[stage, below=6mm of aradev] (diadev) {Diacritizer};
\node[stage, below=6mm of diadev] (ttsdev) {TTS + sync UI};

\draw[myarrow] (bigllm) -- (judgebuild);
\draw[myarrow] (judgebuild) -- (trainloop);
\draw[myarrow] (aradev) -- (diadev);
\draw[myarrow] (diadev) -- (ttsdev);

\draw[myarrow, ultra thick, teal!40!black] (trainloop.east) to[out=0,in=180]
  node[lbl, above, text width=3.2cm, align=center] {one-time export\\quantized weights, offline}
  (aradev.west);
\end{tikzpicture}}
\caption{Deployment boundary. Only quantized AraT5 (and the diacritizer/TTS)
run at inference time on-device; the big LLM and the DSPy judge are strictly
development-time tools whose only artifact that reaches the device is the
trained model weights.}
\label{fig:deployment}
\end{figure}

\label{sec:deployment-note}
\begin{noteb}
\textbf{Pitch framing implication.} The ``on-device AI'' claim is accurate and
defensible \emph{for the simplification model itself} (Stage~1, the intellectual
core). If Stage~3 ends up using a cloud TTS for quality/speed reasons within the
3-week window, that should be stated plainly as a scoped trade-off (``the
reading-support layer currently uses a cloud TTS API; the core NLP model is
on-device''), not glossed over --- judges rewarding ``Pure AI'' clarity will
likely respond better to an honest scoping statement than to an overstated claim
that doesn't survive a follow-up question.
\end{noteb}

\newpage
% =====================================================================
\section{Overall SWOT Analysis}
% =====================================================================

\begin{figure}[h!]
\centering
\begin{tikzpicture}
\def\w{7.6cm}
\def\h{5.6cm}

\node[rectangle, draw=prosgreen, thick, fill=prosgreen!8, minimum width=\w, minimum height=\h,
  align=left, text width=\w-1cm, font=\small, anchor=north east] (S) at (0,0) {
  \textbf{\large\textcolor{prosgreen}{Strengths}}\\[3pt]
  $\bullet$ Open, large, license-unrestricted corpus (BAREC) plus
  confirmed-negligible generation cost --- data and budget risk are both low.\\
  $\bullet$ Every architecture choice has a citable precedent (AraT5-MSAizer).\\
  $\bullet$ Faithfulness-first design matches a vulnerable-user product honestly.\\
  $\bullet$ Rigorous, reproducible QA methodology (DSPy) fits SIC's evident
  preference for demonstrable rigor.\\
  $\bullet$ On-device claim is genuinely defensible for the core model.
};

\node[rectangle, draw=consred, thick, fill=consred!8, minimum width=\w, minimum height=\h,
  align=left, text width=\w-1cm, font=\small, anchor=north west] (W) at (0,0) {
  \textbf{\large\textcolor{consred}{Weaknesses}}\\[3pt]
  $\bullet$ No field-validation partner or real interview subject yet ---
  unlike the team's other Al-Noor-backed ideas.\\
  $\bullet$ Three chained models (simplify, diacritize, judge) is real scope for
  a 3-week sprint.\\
  $\bullet$ No free ground-truth simplification pairs (BAREC has none) ---
  every training pair depends on LLM-generation quality.\\
  $\bullet$ Judge reliability on Arabic semantic equivalence is unverified.\\
  $\bullet$ Mature, timed, on-device Arabic TTS with word timing is a genuinely
  unsolved sub-problem.
};

\node[rectangle, draw=blue!50!black, thick, fill=blue!6, minimum width=\w, minimum height=\h,
  align=left, text width=\w-1cm, font=\small, anchor=south east] (O) at (0,-0.05) {
  \textbf{\large\textcolor{blue!50!black}{Opportunities}}\\[3pt]
  $\bullet$ No deployed Arabic dyslexia software found --- credible first-mover
  framing if it ships.\\
  $\bullet$ Strong accessibility / humanitarian fit with SIC's apparent judging
  preferences (cf.\ past 2nd-place Speech Emotion Synthesis project).\\
  $\bullet$ Week~0 is specifically available to close the partner-org gap before
  the clock starts.\\
  $\bullet$ Runtime faithfulness gate doubles as a genuine, demoable safety
  feature, not just internal QA.
};

\node[rectangle, draw=orange!80!black, thick, fill=orange!8, minimum width=\w, minimum height=\h,
  align=left, text width=\w-1cm, font=\small, anchor=south west] (T) at (0,-0.05) {
  \textbf{\large\textcolor{orange!80!black}{Threats}}\\[3pt]
  $\bullet$ If no partner materializes in week~0, the idea competes at a
  credibility disadvantage against the team's other, interview-backed ideas.\\
  $\bullet$ Three-week timeline is tight for two models plus a judge plus a
  reading UI.\\
  $\bullet$ A competing internal idea could be judged as lower-risk and get
  chosen instead, independent of this idea's own merits.\\
  $\bullet$ BAREC's 19 levels are calibrated for general Arabic readability,
  not dyslexia-specific criteria --- a domain-fit risk that carries over
  regardless of corpus choice.
};
\end{tikzpicture}
\caption{SWOT summary of the project idea as currently scoped.}
\label{fig:swot}
\end{figure}

\newpage
% =====================================================================
\section{Risk Register}
% =====================================================================

\begin{figure}[h!]
\centering
\begin{tikzpicture}[scale=1.0]
\def\n{3}
\foreach \i in {0,...,2} {
  \foreach \j in {0,...,2} {
    \pgfmathsetmacro{\score}{\i+\j}
    \pgfmathsetmacro{\colpct}{25+\score*15}
    \node[rectangle, draw=black!40, minimum width=2.6cm, minimum height=1.9cm,
      fill=red!\colpct!green!60] at (\i*2.6,\j*1.9) {};
  }
}
% axis labels
\node[font=\small] at (-1.6,0) {Low};
\node[font=\small] at (-1.6,1.9) {Medium};
\node[font=\small] at (-1.6,3.8) {High};
\node[rotate=90, font=\bfseries] at (-2.6,1.9) {Impact};
\node[font=\small] at (0,-1.3) {Low};
\node[font=\small] at (2.6,-1.3) {Medium};
\node[font=\small] at (5.2,-1.3) {High};
\node[font=\bfseries] at (2.6,-2.0) {Likelihood};

% risk placements (x = likelihood col 0/1/2, y = impact row 0/1/2)
\node[align=center, font=\scriptsize\bfseries] at (2.6,3.8) {R1, R4};
\node[align=center, font=\scriptsize\bfseries] at (2.6,1.9) {R2, R6};
\node[align=center, font=\scriptsize\bfseries] at (5.2,0) {R5};
\node[align=center, font=\scriptsize\bfseries] at (0,1.9) {R3};
\end{tikzpicture}
\caption{Qualitative likelihood/impact placement for the risks in
Table~\ref{tab:risks}. Placements are judgment calls for team discussion, not a
formal quantitative assessment.}
\label{fig:riskmatrix}
\end{figure}

\begin{center}
\small
\begin{tabular}{@{}p{0.6cm}p{4.4cm}p{3.6cm}p{5.6cm}@{}}
\toprule
\textbf{ID} & \textbf{Risk} & \textbf{Likelihood / Impact} & \textbf{Mitigation} \\
\midrule
R1 & No field-validation partner found during week~0 & Medium / High &
Start outreach immediately (schools, learning-disability orgs, individual
contacts) --- a potential individual validator (someone with dyslexia) has
already been identified as a lead; also consider testing with children, not
just adult readers, since dyslexia support tools are frequently used in
school settings. Worst case, recruit informal testers (any dyslexic
Arabic-literate reader) to at least get qualitative signal before pitching. \\
\addlinespace
R2 & DSPy judge does not reliably agree with human judgment on Arabic semantic
equivalence & Medium / Medium &
Underway: MIPROv2-compile the judge against the growing hand-labeled set
(Section~\ref{sec:qc}) rather than trust it zero-shot; a cheaper
BERTScore-based judge was piloted and rejected for the same reason
($r=-0.014$ with human labels). \\
\addlinespace
R3 & LLM-generated pairs pass the equivalence check but aren't actually
simpler by BAREC's own scale & Low / Medium &
Require a measurable BAREC-calibrated readability-level drop in the judge's
combined check, not equivalence alone; spot-check a sample by hand. \\
\addlinespace
R4 & 3-week timeline insufficient to finish all three stages plus QA
convincingly & Medium / High &
Treat the runtime faithfulness gate as an explicit stretch goal; the demoable
core is Stages~1--3 on typed input with build-time (not runtime) QA only. \\
\addlinespace
R5 & No single on-device Arabic TTS option satisfies quality, commercial
license, and native word timing simultaneously (Section~\ref{sec:tts-arch})
& High / Low--Medium &
Pick per the three-way trade-off in Section~\ref{sec:tts-arch}:
MMS-TTS-Ara or Piper for native timing if the non-commercial license is
acceptable for a capstone demo, else SILMA TTS plus a bolted-on
forced-alignment step; state whichever trade-off is chosen honestly in the
pitch (Section~\ref{sec:deployment-note} framing). If no adequately open
option is found in time, \textbf{drop Stage~3 entirely} rather than ship a
compromised implementation --- a text-only output is still a complete,
demoable product (Section~\ref{sec:tts}). \\
\addlinespace
R6 & Diacritizer trained on classical/religious text underperforms on
simplified contemporary sentences & Medium / Medium &
Validate diacritizer output specifically on Stage~1's simplified sentences
during week~0/1, not just on Tashkeela's native domain; fine-tune further if
needed. \\
\bottomrule
\end{tabular}
\captionof{table}{Risk register corresponding to Figure~\ref{fig:riskmatrix}.}
\label{tab:risks}
\end{center}

\newpage
% =====================================================================
\section{Fit Against the Team's Other Candidate Ideas}
% =====================================================================

\begin{center}
\small
\begin{tabular}{@{}p{3.6cm}p{2.6cm}p{2.6cm}p{2.6cm}p{2.6cm}@{}}
\toprule
& \textbf{This idea}\newline(Arabic dyslexia assistant) & \textbf{Waste-management app} & \textbf{Accessible-camera guidance} & \textbf{RAG curriculum assistant} \\
\midrule
Working prototype exists & No & Yes & Unclear & Unclear \\
Real interview / partner validation & \textbf{No (gap)} & Unclear & Yes (Al-Noor) & Yes (Al-Noor) \\
Public training data available & Yes (BAREC, Tashkeela; pairs LLM-generated) & Unclear & Unclear & Depends on curriculum corpus \\
``Pure AI'' framing strength & Strong (core NLP task) & Weaker (more app/IoT-flavored) & Moderate & Strong \\
On-device story & Strong (small seq2seq) & Depends & Depends (vision on-device is plausible) & Weak (RAG usually needs a larger LM) \\
Accessibility / humanitarian fit & Strong & Weak & Strong & Strong \\
\bottomrule
\end{tabular}
\end{center}

\begin{noteb}
\textbf{Reading this table honestly.} This idea is competitive on data
availability, architectural rigor, and on-device story --- arguably the
strongest of the four on those axes. Its one clear disadvantage against the
Al-Noor-backed ideas is validation: it is the only one without a real interview
subject or partner organization behind it. That single gap is the highest-leverage
thing to close in week~0; closing it would remove the only structural
disadvantage this idea has relative to its siblings.
\end{noteb}

\newpage
% =====================================================================
\section{Execution Timeline}
% =====================================================================

\begin{figure}[h!]
\centering
\resizebox{\textwidth}{!}{%
\begin{ganttchart}[
    hgrid, vgrid,
    x unit=0.62cm,
    y unit title=0.8cm,
    y unit chart=0.7cm,
    title height=1,
    bar height=0.6,
    group height=0.6,
    bar/.append style={fill=stagecolor!80!black, fill=stagecolor},
    group/.append style={fill=blockcolor!80!black, fill=blockcolor},
    milestone/.append style={fill=consred, fill=consred!70}
  ]{1}{20}
  \gantttitle{Week 0}{5}
  \gantttitle{Week 1}{5}
  \gantttitle{Week 2}{5}
  \gantttitle{Week 3}{5} \\
  \gantttitlelist{1,...,20}{1} \\
  \ganttgroup{Validation \& setup}{1}{5} \\
  \ganttbar{Partner/interview outreach}{1}{5} \\
  \ganttbar{Hand-label DSPy val.\ set}{2}{5} \\
  \ganttbar{Pilot BAREC$\to$Gemini generation}{1}{2} \\
  \ganttgroup{Stage 1: Simplify}{4}{12} \\
  \ganttbar{Synthetic data gen (if needed)}{4}{7} \\
  \ganttbar{DSPy judge dev \& pilot}{4}{8} \\
  \ganttbar{AraT5 fine-tuning}{7}{12} \\
  \ganttgroup{Stage 2: Diacritize}{9}{13} \\
  \ganttbar{Integrate \& validate on simplified text}{9}{13} \\
  \ganttgroup{Stage 3: Read aloud (optional)}{6}{15} \\
  \ganttbar{TTS engine selection \& integration}{6}{10} \\
  \ganttbar{Word-sync highlight UI}{10}{15} \\
  \ganttgroup{Integration \& polish}{13}{19} \\
  \ganttbar{End-to-end pipeline wiring}{13}{16} \\
  \ganttbar{Runtime gate (stretch)}{16}{18} \\
  \ganttbar{Testing \& bug-fix buffer}{16}{19} \\
  \ganttmilestone{Pitch / submission prep}{20}
\end{ganttchart}}
\caption{Indicative week-0 plus 3-week sprint schedule (days 1--20, weekdays
only). OCR is out of scope, which frees the days it would have used. The
runtime gate is the first thing cut if the schedule slips; Stage~3 as a
whole is the second, and its optionality (Section~\ref{sec:tts}) means those
days~6--15 fall back to extra buffer for Stages~1--2 and the QC pipeline if
the TTS trade-off in Section~\ref{sec:tts-arch} doesn't resolve in time.}
\label{fig:gantt}
\end{figure}

\newpage
% =====================================================================
\section{Open Decisions and Next Steps}
\label{sec:open-items}
% =====================================================================

\begin{enumerate}[leftmargin=1.6em]
  \item \textbf{Find a real field-validation partner} (a school, a
  learning-disabilities org, or individual dyslexic Arabic readers) --- closes
  the single biggest credibility gap versus the team's other, Al-Noor-backed
  ideas. A potential individual validator (someone with dyslexia) is already
  a lead; testing with children, not just adults, is also worth pursuing
  given dyslexia support tools are frequently used in school settings.
  Highest priority for week~0.
  \item \textbf{Re-evaluate the equivalence validator's batched signature
  against real 5-candidate production data} (Section~\ref{sec:qc}) --- the
  current MIPROv2 calibration is still a batch-of-1 proxy evaluation, and
  the flat 64.0\% thresholded accuracy (validator scoring every held-out
  pair as equivalent) needs to be confirmed as an artifact of that proxy,
  or fixed if it is not, once the 10\% pilot's real batched judgments
  accumulate enough annotated negatives.
  \item \textbf{Continue annotation}, specifically level corrections
  (Section~\ref{sec:annotation}), which remain thinner than equivalence
  corrections; these no longer feed \texttt{ClassifyReadabilityLevel}
  (retired in favor of the supervised MARBERT classifier,
  Section~\ref{sec:next-phase}) but still matter for spot-checking the
  local classifier's real-world error pattern beyond its held-out benchmark
  numbers.
  \item \textbf{Decide whether other BAREC source categories need the same
  identity-pair carve-out as scripture} (Section~\ref{sec:scripture}) ---
  e.g.\ legal or official text --- before generating at full hard-band scale,
  rather than assuming scripture was the only case worth a dedicated rule.
  \item \textbf{Resolve the TTS three-way trade-off, or drop Stage~3}
  (Section~\ref{sec:tts-arch}): MMS-TTS-Ara or Piper for native word timing
  if the CC-BY-NC-4.0 / per-voice license is acceptable, SILMA TTS plus a
  bolted-on forced-alignment step if voice quality and commercial licensing
  matter more --- or, if no adequately open option is found in time, cut
  Stage~3 entirely and ship the text-only pipeline (Section~\ref{sec:tts}).
\end{enumerate}

\begin{noteb}
\textbf{Resolved since the last revision.} (1)~SAMER's license (signed
agreement, no redistribution, non-commercial only) made it a poor fit and
it has been replaced with BAREC (CC~BY-SA~4.0, no agreement required,
larger and more diverse --- Section~\ref{sec:why-barec}); (2)~the API cost
of bulk data generation has been confirmed negligible
(Section~\ref{sec:model-choice}); (3)~the data-splitting design is
concrete and validated --- level~5 dropped (Section~\ref{sec:level5}),
easy/hard/scripture three-way routing at sampling time
(Section~\ref{sec:scripture}), a 2\% pilot actually run (1,350 sentences:
708 easy / 606 hard / 36 scripture); (4)~the DSPy validation set (A3) is
no longer hand-constructed from scratch --- real generated pairs are
annotated directly through a purpose-built review tool
(Section~\ref{sec:annotation}), currently 102 equivalence annotations and
12 level corrections; (5)~the equivalence validator is optimizer-compiled
against real ground truth with measured production accuracy, not run
zero-shot (Section~\ref{sec:qc}); (6)~pipeline resilience --- incremental
checkpointing and per-call timeouts --- was added after a real run-loss
incident (Section~\ref{sec:resilience}); (7)~the level check was moved off
DSPy/DeepSeek entirely onto a fine-tuned, quantized local MARBERT
classifier (86.8--87.5\% band accuracy, up from 79.7\%), and binary
pass/fail generation was replaced with generate-5-then-rerank, both
previously listed as not-yet-built (Section~\ref{sec:next-phase}); (8)~a
cheaper BERTScore-based alternative to the LLM equivalence judge was
piloted and rejected on real annotated data ($r=-0.014$ with human
labels), closing out risk~R2's premise that a non-LLM judge might suffice
(Section~\ref{sec:qc}); (9)~a 10\% stratified pilot sample is in progress
(up from the initial 2\%), annotated through the same review tool.
\end{noteb}

\vspace{0.4cm}
\begin{noteb}
\textbf{Bottom line for the team discussion.} The technical design here is
unusually well-grounded for a capstone idea --- real data, a defensible
architecture choice, and a rigorous QA methodology. The idea's fate mostly hinges
on one non-technical variable: whether a real validation partner can be found in
week~0. Everything else in this document is buildable in three weeks if that one
piece falls into place early.
\end{noteb}

\newpage
% =====================================================================
\section*{References}
\addcontentsline{toc}{section}{References}
% =====================================================================
\begin{enumerate}[leftmargin=1.6em]
\item BAREC: A Large and Balanced Corpus for Fine-grained Arabic Readability
Assessment. arXiv:2502.13520 (ACL 2025 Findings). Replaces SAMER as the
primary data source (Section~\ref{sec:why-barec}).
\item Alhafni, B., et al. ``The SAMER Arabic Text Simplification Corpus.''
arXiv:2404.18615 (LREC-COLING 2024). Considered but not used here --- restrictive
license, see Section~\ref{sec:why-barec}.
\item Zerrouki, T. and Balla, A. ``Tashkeela: Novel corpus of Arabic vocalized
texts, data for auto-diacritization systems.'' \emph{Data in Brief}, 2017.
\item Nagoudi, E.M.B., Elmadany, A., and Abdul-Mageed, M. ``AraT5: Text-to-Text
Transformers for Arabic Language Generation.'' \emph{Proceedings of ACL}, 2022.
\item AraT5-MSAizer. ACL Anthology \texttt{2024.osact-1.16}.
\item ``CATT: Character-based Arabic Tashkeel Transformer.'' arXiv:2407.03236
(ACL ArabicNLP 2024). Recommended diacritization model, Section~\ref{sec:diacritization-arch}.
\item Al-Rfooh, B. and Abandah, G. ``Fine-Tashkeel: Fine-Tuning Byte-Level
Models for Accurate Arabic Text Diacritization.'' arXiv:2303.14588, 2023.
\item Xue, L., et al. ``ByT5: Towards a Token-Free Future with Pre-trained
Byte-to-Byte Models.'' \emph{Transactions of the ACL}, 2022.
\item Khattab, O., Singhvi, A., Maheshwari, P., et al. ``DSPy: Compiling
Declarative Language Model Calls into Self-Improving Pipelines.''
arXiv:2310.03714, 2023.
\item Kim, J., Kong, J., and Son, J. ``Conditional Variational Autoencoder
with Adversarial Learning for End-to-End Text-to-Speech.'' (VITS)
\emph{Proceedings of ICML}, 2021. Architecture behind MMS-TTS and Piper,
Section~\ref{sec:tts-arch}.
\item ``F5-TTS: A Fairytaler that Fakes Fluent and Faithful Speech with Flow
Matching.'' arXiv:2410.06885, 2024. Architecture behind SILMA TTS,
Section~\ref{sec:tts-arch}.
\item ``ArTST: Arabic Text and Speech Transformer.'' arXiv:2310.16621, 2023
(MBZUAI).
\item Related prior work on automatic correction of Arabic dyslexic text, and on
dyslexia-friendly Arabic typeface design (cite specific papers in the final pitch
deck --- referenced here as context confirming the gap, not claiming novelty over
this groundwork).
\end{enumerate}

\subsection*{Vendor and Tooling References}
\begin{itemize}[leftmargin=1.4em]
\item Gemini API pricing (Section~\ref{sec:model-choice}), via
\href{https://developer.puter.com/tutorials/gemini-api-pricing/}{developer.puter.com}.
\item OpenAI API pricing (Section~\ref{sec:model-choice}), via
\href{https://www.cloudzero.com/blog/openai-pricing/}{cloudzero.com}.
\item DeepSeek API pricing, via
\href{https://benchlm.ai/deepseek/api-pricing}{benchlm.ai}.
\item DeepSeek-V4.1-Flash announcement, via
\href{https://www.deepseek.com/en/news/deepseek-v4-1-flash/}{deepseek.com}.
\item SAMER corpus license terms, via
\href{https://camel.abudhabi.nyu.edu/samer-simplification-corpus/}{camel.abudhabi.nyu.edu}.
\item CATT model checkpoints, license, and training code, via
\href{https://github.com/abjadai/catt}{github.com/abjadai/catt}.
\item SILMA TTS v1 model card and details, via
\href{https://huggingface.co/blog/silma-ai/opensource-arabic-english-text-to-speech-model}{huggingface.co/blog/silma-ai}.
\item MMS-TTS-Ara model card (architecture, license, params), via
\href{https://huggingface.co/facebook/mms-tts-ara}{huggingface.co/facebook/mms-tts-ara}.
\item Piper TTS voices and Arabic support, via
\href{https://github.com/rhasspy/piper}{github.com/rhasspy/piper}.
\end{itemize}

\end{document}
